% Part: lambda-calculus
% Chapter: introduction
% Section: computable-lambda

\documentclass[../../../include/open-logic-section]{subfiles}

\begin{document}

\olfileid{lam}{int}{lrp}
\olsection{محاسبه کېدونکې تابعې \usetoken{S}{lambda definable} دي}

\begin{thm}
\ollabel{thm:computable-lambda}
هره محاسبه کېدونکې جزوي تابع !!{lambda definable} ده۔
\end{thm}

\begin{proof}
بايد وښيو چې هره جزوي محاسبه کېدونکې تابع~$f$ يو لامبډا ترم~$F$
!!{lambda defined} کړې ده۔ د کلېنې د عادي بڼې د قضيې له مخې دومره
بس ده چې وښيو هره بنسټيزه بازګشتي تابع يو لامبډا ترم
!!{lambda defined} کړې ده، او بيا دا چې !!{lambda definable} تابعې د مناسب
ترکيبونو او نامحدود لټون تر عملونو لاندې تړلې دي۔ د دې د ښودلو
دپاره چې هره بنسټيزه بازګشتي تابع يو لامبډا ترم !!{lambda defined}
کړې ده، دومره بس ده چې لومړنۍ تابعې !!{lambda definable} وي او
هغه جزوي تابعې چې !!{lambda definable} وي، د ترکيب، بنسټيز بازګښت
او نامحدود لټون تر عملونو لاندې تړلې وي۔
\end{proof}

د ثبوت پاتې برخه د لا آسانه لوستلو دپاره به يوه ډېره دوديزه نښه
وکاروو۔ مثلاً، د $Mxyz$ پر ځاے $M(x, y, z)$ ليکو۔ که څۀ هم دا
نښه څو آرګومېنټونه راښيي، په ياد ولرئ چې د بې‌ټايپه لامبډا حساب
ترمونه ځانته ټاکلې ځاييزه کچه نۀ لري۔ نو د هماغه ترم~$M$ دپاره د
$M(x,y)$ او $M(x,y,z,w)$ ليکل هم معنا لري۔ خو دا نښه ښيي چې موږ
$M$ د درې متغيرونو د تابعې په توګه ګورو، او د تعريفونو تر شا موخه
روښانوي۔ همداسې به د ``$M$ په $M =
\lambd[x][\lambd[y][\lambd[z][\dots]]]$ تعريف کړئ'' پر ځاے
وايو ``$M$ په $M(x,y,z) = \dots$ تعريف کړئ''۔

\end{document}
